\documentclass[10pt,a4paper]{amsart}
\usepackage{amsmath,amssymb,amsthm,graphicx}
\usepackage{verbatim}
\newenvironment{reference}{\comment}{\endcomment}
%\newenvironment{reference}{}{}
\newenvironment{slogan}{\comment}{\endcomment}
\newenvironment{history}{\comment}{\endcomment}
% For commutative diagrams we use Xy-pic
\usepackage[all]{xy}

% We use 2cell for 2-commutative diagrams.
\xyoption{2cell}
% Theorem environments.
%
\theoremstyle{plain}
\newtheorem{theorem}[subsection]{Theorem}
\newtheorem{proposition}[subsection]{Proposition}
\newtheorem{lemma}[subsection]{Lemma}

\theoremstyle{definition}
\newtheorem{definition}[subsection]{Definition}
\newtheorem{example}[subsection]{Example}
\newtheorem{exercise}[subsection]{Exercise}
\newtheorem{situation}[subsection]{Situation}

\theoremstyle{remark}
\newtheorem{remark}[subsection]{Remark}
\newtheorem{remarks}[subsection]{Remarks}

\numberwithin{equation}{subsection}

% Macros
%
\def\lim{\mathop{\mathrm{lim}}\nolimits}
\def\colim{\mathop{\mathrm{colim}}\nolimits}
\def\Spec{\mathop{\mathrm{Spec}}}
\def\Hom{\mathop{\mathrm{Hom}}\nolimits}
\def\Ext{\mathop{\mathrm{Ext}}\nolimits}
\def\SheafHom{\mathop{\mathcal{H}\!\mathit{om}}\nolimits}
\def\SheafExt{\mathop{\mathcal{E}\!\mathit{xt}}\nolimits}
\def\Sch{\mathit{Sch}}
\def\Mor{\mathop{\mathrm{Mor}}\nolimits}
\def\Ob{\mathop{\mathrm{Ob}}\nolimits}
\def\Sh{\mathop{\mathit{Sh}}\nolimits}
\def\NL{\mathop{N\!L}\nolimits}
\def\CH{\mathop{\mathrm{CH}}\nolimits}
\def\proetale{{pro\text{-}\acute{e}tale}}
\def\etale{{\acute{e}tale}}
\def\QCoh{\mathit{QCoh}}
\def\Ker{\mathop{\mathrm{Ker}}}
\def\Im{\mathop{\mathrm{Im}}}
\def\Coker{\mathop{\mathrm{Coker}}}
\def\Coim{\mathop{\mathrm{Coim}}}

% Boxtimes
%
\DeclareMathSymbol{\boxtimes}{\mathbin}{AMSa}{"02}

%
% Macros for moduli stacks/spaces
%
\def\QCohstack{\mathcal{QC}\!\mathit{oh}}
\def\Cohstack{\mathcal{C}\!\mathit{oh}}
\def\Spacesstack{\mathcal{S}\!\mathit{paces}}
\def\Quotfunctor{\mathrm{Quot}}
\def\Hilbfunctor{\mathrm{Hilb}}
\def\Curvesstack{\mathcal{C}\!\mathit{urves}}
\def\Polarizedstack{\mathcal{P}\!\mathit{olarized}}
\def\Complexesstack{\mathcal{C}\!\mathit{omplexes}}
% \Pic is the operator that assigns to X its picard group, usage \Pic(X)
% \Picardstack_{X/B} denotes the Picard stack of X over B
% \Picardfunctor_{X/B} denotes the Picard functor of X over B
\def\Pic{\mathop{\mathrm{Pic}}\nolimits}
\def\Picardstack{\mathcal{P}\!\mathit{ic}}
\def\Picardfunctor{\mathrm{Pic}}
\def\Deformationcategory{\mathcal{D}\!\mathit{ef}}
\usepackage{cleveref}
\crefname{originallemma}{original lemma}{original lemmas}
\Crefname{originallemma}{Original lemma}{Original lemmas}
\makeatletter
\newlabel{lemma-base-change-separated}{{lemma-base-change-separated}{source}{}{sourcepointer-1}{}}
\newlabel{lemma-base-change-separated@cref}{{[originallemma][0][]lemma-base-change-separated}{[1][0][]source}{}{}{}}
\makeatother
\crefname{originallemma}{original lemma}{original lemmas}
\Crefname{originallemma}{Original lemma}{Original lemmas}
\makeatletter
\newlabel{lemma-separated-over-separated}{{lemma-separated-over-separated}{source}{}{sourcepointer-2}{}}
\newlabel{lemma-separated-over-separated@cref}{{[originallemma][0][]lemma-separated-over-separated}{[1][0][]source}{}{}{}}
\makeatother
\crefname{originallemma}{original lemma}{original lemmas}
\Crefname{originallemma}{Original lemma}{Original lemmas}
\makeatletter
\newlabel{lemma-trivial-implications}{{lemma-trivial-implications}{source}{}{sourcepointer-3}{}}
\newlabel{lemma-trivial-implications@cref}{{[originallemma][0][]lemma-trivial-implications}{[1][0][]source}{}{}{}}
\makeatother
\crefname{originaltheorem}{original theorem}{original theorems}
\Crefname{originaltheorem}{Original theorem}{Original theorems}
\makeatletter
\newlabel{theorem-DM}{{theorem-DM}{source}{}{sourcepointer-4}{}}
\newlabel{theorem-DM@cref}{{[originaltheorem][0][]theorem-DM}{[1][0][]source}{}{}{}}
\makeatother
\title{A base change of a DM morphism along a DM stack or an algebraic space admits a surjective etale cover by a scheme}
\author{The Stacks Project Authors}
\date{stacks-project:stacks-morphisms.tex:local-snapshot-20261012; GNU FDL 1.2 or later}
\begin{document}
\maketitle
\noindent\textit{Editorial note.} All delivered mathematical text is the author's own, copied from the pinned Stacks Project chapter source sources/stacks-algebraic-stacks-morphisms/pinned-original-source.tex; no statement or proof step was rewritten, reordered or removed.\par\medskip
\noindent\textit{Editorial note.} Preamble provenance: the Stacks Project keeps one single shared preamble for every chapter; this card takes it directly from the staged sources/stacks-algebraic-stacks-morphisms/preamble.tex as three verbatim spans (lines 8-12, 14-18 and 154-217), so no derived file was written and nothing was added or redefined. The wrapper class and the base packages are supplied by the assembler.\par\medskip
\noindent\textit{Editorial note.} Bibliography provenance: neither the counted statement nor its proof contains a citation (this lane rescans the whole span before writing the card), so the card carries no bibliography and no entry was reconstructed.\par\medskip
\noindent\textit{Editorial note.} Selection provenance: this card was opened under the file-based control contract (CONTROL.md), with the mechanical pre-card gate run on its statement and proof interval before the card was written.\par\medskip
\noindent\textit{Allowed standard prerequisites (editorial).}\begin{itemize}
\item The base change permanence of DM-ness. Invoked at the start of the counted proof (source lines 4762-4764). Makes the base-changed morphism DM.
\item The statement that a morphism between stacks that is separated over a base and has DM target is DM. Invoked inside the counted proof (source lines 4765-4767). Makes the base-changed stack DM.
\item The characterization of DM stacks through the existence of a surjective etale cover by a scheme. Invoked at the end of the counted proof (source lines 4768-4770). Produces the scheme and the surjective etale morphism asserted.
\end{itemize}\medskip
\ifdefined\hypertarget\hypertarget{sourcepointer-1}{}\fi
\noindent\textit{Source reference.} Original lemma lemma-base-change-separated: sources/stacks-algebraic-stacks-morphisms/pinned-original-source.tex, source lines 703--703. Navigation pointer only: the counted text refers to the Stacks item lemma-base-change-separated, which is an existing labeled result of this same chapter and is not part of this card. The pointer carries the target's own label as its number and its definition line as its span; its content is not reproduced here.\par\medskip
\ifdefined\hypertarget\hypertarget{sourcepointer-2}{}\fi
\noindent\textit{Source reference.} Original lemma lemma-separated-over-separated: sources/stacks-algebraic-stacks-morphisms/pinned-original-source.tex, source lines 888--888. Navigation pointer only: the counted text refers to the Stacks item lemma-separated-over-separated, which is an existing labeled result of this same chapter and is not part of this card. The pointer carries the target's own label as its number and its definition line as its span; its content is not reproduced here.\par\medskip
\ifdefined\hypertarget\hypertarget{sourcepointer-3}{}\fi
\noindent\textit{Source reference.} Original lemma lemma-trivial-implications: sources/stacks-algebraic-stacks-morphisms/pinned-original-source.tex, source lines 682--682. Navigation pointer only: the counted text refers to the Stacks item lemma-trivial-implications, which is an existing labeled result of this same chapter and is not part of this card. The pointer carries the target's own label as its number and its definition line as its span; its content is not reproduced here.\par\medskip
\ifdefined\hypertarget\hypertarget{sourcepointer-4}{}\fi
\noindent\textit{Source reference.} Original theorem theorem-DM: sources/stacks-algebraic-stacks-morphisms/pinned-original-source.tex, source lines 4589--4589. Navigation pointer only: the counted text refers to the Stacks item theorem-DM, which is an existing labeled result of this same chapter and is not part of this card. The pointer carries the target's own label as its number and its definition line as its span; its content is not reproduced here.\par\medskip
\begin{lemma}
\label{lemma-DM}
Let $f : \mathcal{X} \to \mathcal{Y}$ be a DM morphism of algebraic
stacks. Then
\begin{enumerate}
\item For every DM algebraic stack $\mathcal{Z}$ and morphism
$\mathcal{Z} \to \mathcal{Y}$ there exists a scheme and
a surjective \'etale morphism
$U \to \mathcal{X} \times_\mathcal{Y} \mathcal{Z}$.
\item For every algebraic space $Z$ and morphism
$Z \to \mathcal{Y}$ there exists a scheme and
a surjective \'etale morphism
$U \to \mathcal{X} \times_\mathcal{Y} Z$.
\end{enumerate}
\end{lemma}

\begin{proof}
Proof of (1). As $f$ is DM we see that the base change
$\mathcal{X} \times_\mathcal{Y} \mathcal{Z} \to \mathcal{Z}$ is DM
by Lemma \ref{lemma-base-change-separated}.
Since $\mathcal{Z}$ is DM this implies that
$\mathcal{X} \times_\mathcal{Y} \mathcal{Z}$ is DM
by Lemma \ref{lemma-separated-over-separated}. Hence there exists a
scheme $U$ and a surjective \'etale morphism
$U \to \mathcal{X} \times_\mathcal{Y} \mathcal{Z}$, see
Theorem \ref{theorem-DM}.
Part (2) is a special case of (1) since an algebraic space
(when viewed as an algebraic stack) is DM by
Lemma \ref{lemma-trivial-implications}.
\end{proof}

\end{document}
